\subsection{Постановка задачи распределения вершин графа}

В контексте распределенных графовых баз данных задача распределения вершин (графового партиционирования) формулируется следующим образом. Пусть дан неориентированный граф \(G = (V, E)\), где \(V\) - множество вершин, \(|V| = n\), а \(E\) - множество рёбер, \(|E| = m\). Требуется найти такое разбиение (распределение) \(P = \{S_1, S_2, ..., S_k\}\) множества вершин \(V\) на \(k\) непересекающихся подмножеств (шардов) \(S_i\), что:
\begin{itemize}
    \item \(\bigcup_{i=1}^{k} S_i = V\);
    \item \(S_i \cap S_j = \emptyset\) для всех \(i \neq j\).
\end{itemize}

Качество распределения оценивается по двум основным критериям, которые необходимо оптимизировать:
\begin{itemize}
    \item минимизация разрезов (Edge Cut). Мощность разреза \(\partial e(P)\) определяется как количество рёбер, концы которых оказались в разных шардах. Формально:
    $$
    |\partial e(P)| = \sum_{i=1}^{k} |e(S_i, V \setminus S_i)|
    $$
    Минимизация этого параметра напрямую влияет на снижение сетевого трафика при выполнении распределенных запросов.
    \item балансировка нагрузки (Load Balancing). Неравномерное распределение вершин приводит к перегрузке одних узлов и простою других. Для количественной оценки используется показатель нормализованной максимальной нагрузки \(\rho\), который необходимо минимизировать:
    $$
    \rho = \frac{\max_{i=1}^{k} (|S_i|)}{n/k}
    $$
    Идеально сбалансированное распределение соответствует \(\rho = 1\).
\end{itemize}
Таким образом, целевая задача заключается в поиске такого распределения \(P^*\), которое минимизирует мощность разрезов \(|\partial e(P)|\) при максимально допустимом значении дисбаланса \(\rho \leq \tau\), где \(\tau\) - заданный порог.
